% Directly authored companion to gamma_uniform_grid_theorem.md.
\documentclass[11pt]{article}
\usepackage[T1]{fontenc}
\usepackage{amsmath,amssymb,amsthm}
\usepackage{newtxtext,newtxmath}
\usepackage[letterpaper,margin=0.78in]{geometry}
\usepackage{graphicx,microtype,xcolor,caption,fancyhdr,tabularx}
\usepackage[colorlinks=true,linkcolor=blue!40!black,urlcolor=blue!50!black]{hyperref}
\definecolor{ink}{HTML}{263442}
\DeclareMathOperator{\sech}{sech}
\DeclareMathOperator{\sign}{sign}
\DeclareMathOperator{\diag}{diag}
\DeclareMathOperator{\rank}{rank}
\theoremstyle{definition}
\newtheorem{lemma}{Lemma}
\newtheorem{proposition}[lemma]{Proposition}
\newtheorem{theorem}[lemma]{Theorem}
\newtheorem{corollary}[lemma]{Corollary}
\newcommand{\repo}{https://github.com/samlayton99/precision-mlps/blob/experiment/frozen-gamma-probe}
\newcommand{\evidence}{\repo/results/checkpoint_D_optimizers/expD36_frozen_gamma_probe/full_sweep/refinements/uniform_grid_spectrum}
\graphicspath{{results/checkpoint_D_optimizers/expD36_frozen_gamma_probe/full_sweep/refinements/uniform_grid_spectrum/}}
\hypersetup{pdftitle={How a common tanh slope controls the readout spectrum},
 pdfsubject={Exact periodic gamma law, finite uniform-grid boundary corrections, proofs, and acquisition bounds}}
\setlength{\parindent}{0pt}
\setlength{\parskip}{5pt plus 1pt}
\setlength{\emergencystretch}{2em}
\setlength{\headheight}{14pt}
\captionsetup{font=small,labelfont=bf,justification=raggedright,singlelinecheck=false}
\pagestyle{fancy}
\fancyhf{}
\fancyhead[L]{\small\color{ink}Gamma and the uniform-grid readout spectrum}
\fancyhead[R]{\small Full proof note}
\fancyfoot[C]{\small\thepage}
\renewcommand{\headrulewidth}{0.3pt}

\begin{document}
\thispagestyle{plain}
{\color{ink}\fontsize{23}{26}\selectfont\bfseries How a common tanh slope\\controls the readout spectrum\par}
\vspace{8pt}
Small gamma smooths every feature before the readout sees it. On a periodic,
uniform geometry, this smoothing multiplies each frequency's learning rate
by an explicit factor squared. That statement is exact: the frequencies are
the kernel's independent learning directions. On a finite interval, ordinary
tanh has additional boundary effects. We retain those effects explicitly as
a finite correction, while bounding the remaining error. This gives a
gamma-dependent spectral construction with a controllable error, rather than
a single worst-case inequality whose slack grows unchecked.

There are two different conclusions. The periodic theorem provides a closed
frequency-by-frequency law. The ordinary-tanh theorem provides a diagonalizable
reference plus a signed low-rank correction; predicting its individual
eigenvalues still involves a finite spectral calculation. The correction is
not assumed negligible, and its rank must be reported. Neither uniform spacing
nor this proof implies that every target learns faster as gamma increases on
the finite interval.

The numerical check below retains that correction. At gamma 8, it recovers
the observed slow target energy with a narrow enclosure and predicts the
executed 15.8-million-update acquisition time. The 26th normalized eigenvalue
has a relative enclosure width of about 0.00171\%, compared with the earlier
simplified upper bound being approximately 32,000 times its true value.
These spectral evaluations use a disclosed floating-point allowance; the
primary timing intervals additionally follow from existing independent
interval certificates. The useful finite result is the corrected spectral
theorem, not the bare periodic formula or a rank-shift bound with the
correction discarded.

\textbf{Symbols and normalization.} Eigenvalues are denoted by $\mu$ or $\nu$;
$\lambda$ denotes dimensionless bandwidth, not an eigenvalue.

\begin{center}\small
\begin{tabularx}{\linewidth}{@{}p{0.24\linewidth}X@{}}
\hline
Symbol & Meaning\\\hline
$\gamma>0$ & Common, frozen tanh slope\\
$L$ & Physical interval length; reference period is $2L$\\
$N,h=L/N$ & Number of core centers and their spacing\\
$q,m=qN+1$ & Integer sampling oversampling factor and number of inputs\\
$W$ & Total tanh feature count, including any halo centers\\
$J,K=JJ^*$ & Feature matrix normalized by $\sqrt m$, and raw-readout kernel\\
$p_\gamma,M_\gamma(\omega)$ & Smoothed antiperiodic reference and frequency attenuation\\
$K_0,\widehat K$ & Reference kernel and kernel retaining boundary corrections\\
$B,p$ & Exact columns at each boundary and retained exponential terms\\
$\delta$ & Upper bound on $\|K-\widehat K\|$\\
$S_K(s)$ & Fraction of target energy in eigenvalues at most $s$\\
$\eta,n$ & GD step and update count\\
$\lambda=\gamma h$ & Slope measured relative to center spacing\\\hline
\end{tabularx}
\end{center}

\section{The optimization problem is unchanged}
For physical target samples $y_i^{\rm phys}$, let
$y_i=y_i^{\rm phys}/\sqrt m$. The feature matrix has the bias column
$1/\sqrt m$ and columns
\[
J_{ij}=\frac{\tanh(\gamma(x_i-c_j))}{\sqrt m}.
\]
The loss is $\frac12\|Jw-y\|^2$, equal to the original mean squared residual
divided by two. We optimize the original readout coefficients $w$, with their
ordinary Euclidean metric. Starting from $w_0=0$, the normalized residual is
\[
r_n=(I-\eta K)^n y,\qquad K=JJ^*,\qquad
E(n)=\frac{\|r_n\|}{\|y\|}.
\]
We assume $y\ne0$ and $\eta>0$. In training statements choose $\eta\|K\|\le1$, so each
component decays without alternating sign. These are training residuals on
the specified inputs, not held-out generalization claims.

A change of feature coordinates is permitted only if its metric is retained.
For example, with $S_j=\sum_{\ell\le j}w_\ell$,
\begin{align*}
\sum_{j=1}^{W}w_j\tanh(\gamma(x-c_j))
={}&S_W\tanh(\gamma(x-c_W))\\
&+\sum_{j=1}^{W-1}S_j
[\tanh(\gamma(x-c_j))-\tanh(\gamma(x-c_{j+1}))].
\end{align*}
This identity does not authorize ordinary Euclidean GD in $S$: the original
norm is $\sum_j(S_j-S_{j-1})^2$, with $S_0=0$. The construction below works
directly with $J$, so this metric issue never becomes an approximation.

\section{Gamma is an exact smoothing multiplier}
\begin{lemma}[Smoothing a step]
Define the probability density
\[
\rho_\gamma(t)=\frac\gamma2\sech^2(\gamma t).
\]
Then
\[
(\rho_\gamma*\sign)(x)=\tanh(\gamma x),\qquad
\widehat\rho_\gamma(\omega)=M_\gamma(\omega)
=\frac{z}{\sinh z},\qquad z=\frac{\pi|\omega|}{2\gamma},
\]
with $M_\gamma(0)=1$ by continuity.
\end{lemma}
\begin{proof}
The convolution is twice the distribution function of $\rho_\gamma$ minus
one, giving tanh. For the Fourier transform, substitute $u=e^{2\gamma t}$:
\[
\widehat\rho_\gamma(\omega)
=\int_0^\infty\frac{u^{-i\omega/(2\gamma)}}{(1+u)^2}\,du
=\Gamma(1-ia)\Gamma(1+ia)
=\frac{\pi a}{\sinh(\pi a)},\qquad a=\frac\omega{2\gamma}.
\]
The middle equality is the beta integral and the last is the gamma reflection
identity. Evenness gives the formula. For $z>0$, $z/\sinh z$ decreases because
$\sinh z-z\cosh z<0$; hence attenuation increases with gamma at fixed nonzero
frequency. Also,
\[
M_\gamma(\omega)=\frac{2ze^{-z}}{1-e^{-2z}}.
\]
Thus high frequency relative to gamma is suppressed exponentially.
\end{proof}
This describes the feature construction before any target or optimizer is
chosen.

\section{An exact spectral law under periodic uniform geometry}
Let $s$ be the $2L$-periodic square wave equal to $\sign(x)$ on $(-L,L)$, and
let $p_\gamma=\rho_\gamma*s$, using the periodic extension. Its nonzero
complex Fourier coefficients are
\[
b_k=\frac2{i\pi(2k+1)},\qquad
\omega_k=\frac{(2k+1)\pi}{L},\qquad k\in\mathbb Z,
\]
and
\[
p_\gamma(x)=\sum_{k\in\mathbb Z}
b_kM_\gamma(\omega_k)e^{i\omega_kx}.
\]
For positive gamma these coefficients decay exponentially, giving absolute
convergence. Notice $p_\gamma(x+L)=-p_\gamma(x)$: \emph{antiperiodic} means
that a shift by $L$ flips its sign, while a shift by $2L$ returns it unchanged.

\begin{theorem}[Periodic gamma-to-spectrum law]
With centers uniform over a full period and the same uniform sampling measure,
define the center-average kernel
\[
\mathcal K_\gamma(x,x')
=\frac1{2L}\int_0^{2L}p_\gamma(x-c)p_\gamma(x'-c)\,dc.
\]
On $L^2([0,2L],dx/(2L))$, each odd Fourier wave has eigenvalue
\[
\mu_k(\gamma)=|b_k|^2M_\gamma(\omega_k)^2.
\]
Every nonconstant even Fourier wave has eigenvalue zero. Adding a unit bias
gives the constant wave eigenvalue one without changing the other modes.
If total feature mass is $a$ instead of one, multiply nonbias eigenvalues
by $a$; a raw finite-width readout is not width-normalized without explicitly
introducing that factor.
\end{theorem}
\begin{proof}
Substitute the absolutely convergent Fourier expansions. Integration over
$c$ eliminates unequal frequencies, giving
\[
\mathcal K_\gamma(x,x')
=\sum_k|b_k|^2M_\gamma(\omega_k)^2
e^{i\omega_kx}e^{-i\omega_kx'}.
\]
Orthogonality under the input measure gives the eigenvalues. The square wave
has no even Fourier coefficients. The constant bias is orthogonal to every
nonconstant Fourier wave.
\end{proof}
This is exact, not a norm upper bound. For a single accessible Fourier target,
GD has exact relative residual
\[
E(n)=(1-\eta\mu_k(\gamma))^n.
\]
For $0<\varepsilon<1$ and $0<\eta\mu_k<1$, its acquisition time is exactly
\[
\left\lceil
\frac{\log(1/\varepsilon)}{-\log(1-\eta\mu_k(\gamma))}
\right\rceil.
\]
A cap $\gamma\le\Gamma$ implies
$\mu_k(\gamma)\le|b_k|^2M_\Gamma(\omega_k)^2$ and therefore a necessary
learning time at a fixed stable step. For several specified frequencies,
sum their exact decays. Missing even modes cause an irreducible residual,
not slow convergence. These distinctions matter whenever periodic square
waves serve as the illustrative model.

The remainder transfers the explicit frequency mechanism to the ordinary
finite-interval tanh dictionary. It does not identify that dictionary with
the periodic population model.

\section{The finite uniform core is diagonalizable, with aliases retained}
Assume positive integers $N,q$, $h=L/N$, inputs $x_i=ih/q$ for
$i=0,\ldots,qN$, and core centers $c_j=jh$ for $j=0,\ldots,N-1$.
Additional centers may be included and will be treated explicitly. The
ordinary-tanh matrix includes the bias.

For $J_0$, use $p_\gamma(x_i-c_j)/\sqrt m$ on the first $qN$ rows and
the $N$ core columns. Set its final row, additional columns, and bias column
to zero. This makes $J_0$ the same size as $J$ without changing the
reference's nonzero spectrum.

\begin{proposition}[Exact finite alias blocks]
Reorder the first $qN$ rows as $i=q\ell+s$, with $\ell=0,\ldots,N-1$ and
$s=0,\ldots,q-1$. Define
\[
T_s[\ell,j]=p_\gamma((\ell-j)h+sh/q),\qquad
Q[\ell,k]=N^{-1/2}e^{i(2k+1)\pi\ell/N}.
\]
Then $Q$ is unitary and
\[
T_s=Q\diag(t_{s,0},\ldots,t_{s,N-1})Q^*,
\]
where the exact alias formula is
\[
t_{s,k}=N\sum_{a\in\mathbb Z}
b_{k+aN}M_\gamma(\omega_{k+aN})
e^{i\omega_{k+aN}sh/q}.
\]
Consequently $K_0=J_0J_0^*$ decomposes into $N$ rank-at-most-one
$q\times q$ blocks
\[
\frac1m t_kt_k^*,\qquad t_k=(t_{0,k},\ldots,t_{q-1,k})^T,
\]
and a zero block. Its possible nonzero eigenvalues are
\[
\nu_k(\gamma)=\frac1m\sum_{s=0}^{q-1}|t_{s,k}|^2.
\]
\end{proposition}
\begin{proof}
Antiperiodicity makes $T_s$ skew-circulant: wrapping a center difference
by $N$ grid spacings changes the sign. Substituting a column of $Q$ in $T_sQ$
gives an eigenvector with eigenvalue
\[
\sum_{\ell=0}^{N-1}p_\gamma(\ell h+sh/q)
e^{-i(2k+1)\pi\ell/N}.
\]
Insert the Fourier series. The finite geometric sum equals $N$ precisely
when the Fourier index equals $k$ modulo $N$, and otherwise vanishes.
This gives the alias formula. Applying $Q^*$ independently to each phase
group and grouping equal $k$ gives the stated blocks. The normalization is
$m=qN+1$, even though the reference's last row is zero.
\end{proof}
For $q=1$, aliases are added before taking the squared magnitude. They
cannot be replaced by a sum of squared coefficients. With oversampling,
each block keeps relative phases between sample groups. Gamma is explicit
in every term, but monotonicity of an individual aliased eigenvalue does
not follow from monotonicity of each attenuation factor.

An implementation may evaluate $p_\gamma$ through the image formula below
instead of truncated Fourier modes. These evaluate the same reference;
any truncation in either route belongs in the final error budget.

\section{Ordinary tanh equals the reference plus controlled boundary terms}
\begin{lemma}[Exact image representation]
For real $d$,
\[
p_\gamma(d)=\tanh(\gamma d)+\sum_{\ell=1}^\infty(-1)^\ell
[\tanh(\gamma(d-\ell L))+\tanh(\gamma(d+\ell L))].
\]
This paired series is locally uniformly convergent. For $|d|<L$,
\[
\tanh(\gamma d)-p_\gamma(d)
=2\sum_{r=1}^\infty
\frac{(-1)^{r-1}}{1+e^{-2r\gamma L}}
[e^{-2r\gamma(L-d)}-e^{-2r\gamma(L+d)}].
\]
\end{lemma}
\begin{proof}
Distributionally, the derivative of the periodic square wave is
$2\sum_{\ell\in\mathbb Z}(-1)^\ell\delta_{\ell L}$.
Convolving with $\rho_\gamma$ gives
$\sum_\ell(-1)^\ell\gamma\sech^2(\gamma(d-\ell L))$.
This is also the derivative of the paired image series. Both functions
vanish at $d=0$, so they agree. Exponential decay of paired tails justifies
local convergence and differentiation.

For $a>0$, the geometric expansion is
\[
\tanh(\gamma a)=1-2\sum_{r=1}^\infty(-1)^{r-1}e^{-2r\gamma a}.
\]
In each image pair, $\ell L-d$ and $\ell L+d$ are positive for $|d|<L$.
Substitute their expansions and sum
\[
\sum_{\ell\ge1}(-1)^{\ell+1}e^{-2r\gamma\ell L}
=\frac{e^{-2r\gamma L}}{1+e^{-2r\gamma L}}.
\]
Absolute convergence on compact subsets of $(-L,L)$ permits the interchange.
\end{proof}

\begin{lemma}[Finite-rank boundary approximation and its tail]
Choose an integer $B$ with $1\le B<N/2$ and a nonnegative integer $p$.
Treat core columns $j<B$ and $j\ge N-B$ exactly. On the other core columns
and first $qN$ rows, retain the first $p$ terms of Lemma~4. Define
\[
\epsilon_{B,p}(\gamma)=
\frac{4e^{-2(p+1)\gamma Bh}}{1-e^{-2(p+1)\gamma L}}.
\]
Every remaining unnormalized feature entry has error at most
$\epsilon_{B,p}(\gamma)$.
\end{lemma}
\begin{proof}
A tanh expansion through $p$ terms has absolute remainder at most
$2e^{-2(p+1)\gamma a}$. The finite geometric identity gives this directly:
the exact remainder has an additional denominator
$1+e^{-2\gamma a}\ge1$. Truncate each positive-distance tanh in each image
pair before summing images. Interior columns have $|d|\le L-Bh$, so
\begin{align*}
|\mathrm{remainder}|
&\le2\sum_{\ell=1}^\infty
\left[e^{-2(p+1)\gamma(\ell L-d)}
+e^{-2(p+1)\gamma(\ell L+d)}\right]\\
&\le4\sum_{\ell=1}^\infty
e^{-2(p+1)\gamma((\ell-1)L+Bh)}
=\epsilon_{B,p}(\gamma).
\end{align*}
Summing the retained terms over images gives exactly the first $p$ terms
in Lemma~4. Thus the bound concerns the claimed approximation, not a
different truncation.
\end{proof}
Each retained term separates into a function of $x$ times a function of $c$:
\[
e^{-2r\gamma(L-d)}-e^{-2r\gamma(L+d)}
=e^{-2r\gamma(L-x)}e^{-2r\gamma c}
-e^{-2r\gamma x}e^{-2r\gamma(L-c)}.
\]
All factors are at most one on $x,c\in[0,L]$. This form avoids overflow and
proves the rank bound without changing the parameter metric.

Construct $\widehat J$ by applying the core corrections to $J_0$, restoring
boundary columns exactly, adding every noncore tanh column and bias exactly,
and restoring the last input row exactly. If $e=W-N$ is the number of
noncore tanh columns, then
\[
\widehat J=J_0+UV^*,\qquad
\rank(UV^*)\le r_f:=2B+2p+e+2.
\]
Here $2B$ counts exactly corrected core columns, $2p$ the separable
exponentials, $e+1$ the additional features and bias, and one the last row.
This is an upper bound, not a claim that all components are independent.
Masks restricting exponential terms to interior columns and the first
$qN$ rows preserve their rank. Endpoint corrections are applied last.

At most $N-2B$ columns have approximation error. After normalization,
\[
\|J-\widehat J\|\le\|J-\widehat J\|_F
\le\epsilon_J:=\sqrt{N-2B}\,\epsilon_{B,p}(\gamma).
\]
Every original feature, including bias, has absolute unnormalized value
at most one; hence $\|J\|\le\sqrt{W+1}$ and
\[
\|K-\widehat K\|\le
\delta:=2\sqrt{W+1}\epsilon_J+\epsilon_J^2,
\qquad\widehat K=\widehat J\widehat J^*.
\]
Expand $\widehat J\widehat J^*-JJ^*$ and apply the triangle inequality to
obtain the last bound. For fixed $B>0$ and gamma, it decreases geometrically
to zero as $p$ increases. It bounds analytic truncation, not floating-point
error. Reference evaluation, factorization, and eigensolver errors require
their own accounting before a numerical enclosure can be called certified.

\section{The finite-interval spectral theorem}
\begin{theorem}[Explicit bulk, retained boundary correction, and enclosures]
Under the uniform geometry above,
\[
K=K_0+ZHZ^*+E,\qquad\|E\|\le\delta,
\]
where Proposition~3 gives the bulk spectrum and
\[
Z=[J_0V,U],\qquad
H=\begin{bmatrix}0&I\\I&V^*V\end{bmatrix}.
\]
The correction is signed Hermitian of rank at most $r_b=2r_f$.
It need not be positive semidefinite, although $K$ and
$\widehat K=K_0+ZHZ^*$ are positive semidefinite.

Let $\mu_j,\widehat\mu_j,\nu_j$ be the eigenvalues of $K,\widehat K,K_0$
in decreasing order, including zeros. Then
\[
|\mu_j-\widehat\mu_j|\le\delta.
\]
Where both shifted indices exist,
\[
(\nu_{j+r_b}-\delta)_+\le\mu_j
\le\nu_{j-r_b}+\delta.
\]
A sharper known correction rank may replace $r_b$.
\end{theorem}
\begin{proof}
Expand the feature Gram matrix:
\[
\widehat K-K_0
=J_0VU^*+UV^*J_0^*+U(V^*V)U^*=ZHZ^*.
\]
Its rank is at most the number of columns of $Z$. Lemma~5 bounds the
remaining operator error. The min--max characterization gives Weyl's norm
bound. Intersecting a candidate min--max subspace with the nullspace of the
rank-$r_b$ correction reduces its dimension by at most $r_b$; apply this
once for the upper bound and once for the lower bound to obtain rank
interlacing. Finally use nonnegativity of $K$.
\end{proof}
Index-shift inequalities can be conservative. To keep boundary information,
diagonalize the signed low-rank update instead. In a basis diagonalizing
$K_0$, write $D=\diag(\nu_j)$ and $\overline Z$ for the transformed factors.
Away from the poles $z=\nu_j$, the matrix determinant lemma gives
\[
\det(\widehat K-zI)=\det(D-zI)
\det\!\left(I+H\overline Z^*(D-zI)^{-1}\overline Z\right).
\]
The second determinant has dimension at most $2r_f$. At a pole, interpret
the product by continuity or use deflation; simply skipping poles does not
give a complete eigensolver. Unchanged eigenvectors orthogonal to the
correction must also be retained. This is an exact representation of
$\widehat K$, providing a structured finite problem rather than a short
scalar formula for every eigenvalue.

\textbf{Coefficient-space version and compressed correction.}
When there are fewer features than inputs, use $G=J^*J$. It has exactly
the same positive eigenvalues as $K$; remaining dimensions contribute zeros.
With $G_0=J_0^*J_0$, the corresponding expansion is
\[
\widehat G=\widehat J^*\widehat J=G_0+Z_cB_cZ_c^*,\qquad
Z_c=[J_0^*U,V],\qquad
B_c=\begin{bmatrix}0&I\\I&U^*U\end{bmatrix}.
\]
The feature-error estimate also gives $\|G-\widehat G\|\le\delta$.
Thus coefficient space suffices for eigenvalue enclosures without forming
the larger sample kernel. Target projections still live in sample space;
recover them, for example, from singular vectors of $\widehat J$.
Coefficient eigenvalues alone do not supply these projections.

Diagonalize $\widehat G-G_0$ on its range. Retain correction eigenpairs with
$|s_\ell|>\tau$, giving orthonormal columns $Q_c$ and
$\Sigma=\diag(s_\ell)$. Discarding the remaining terms has norm at most
$\tau$ in exact arithmetic. Therefore
\[
G_{\rm comp}=G_0+Q_c\Sigma Q_c^*,\qquad
\|G-G_{\rm comp}\|\le\delta+\tau.
\]
These correction eigenvalues are signed: truncation by magnitude, rather
than retaining only positive terms, is essential. Numerical compression
error is additional to this exact-arithmetic allowance.

In a basis diagonalizing $G_0$, call its diagonal matrix $D_c$ and the
transformed retained columns $\overline Q_c$. Let $n_-(A)$ count strictly
negative eigenvalues of a Hermitian matrix. At a threshold $t$ away from
the diagonal entries of $D_c$,
\begin{align*}
n_-(G_{\rm comp}-tI)
={}&n_-(D_c-tI)\\
&+n_-\!\left(-\Sigma^{-1}
-\overline Q_c^*(D_c-tI)^{-1}\overline Q_c\right)
-n_-(-\Sigma^{-1}).
\end{align*}
\begin{proof}[Proof of the count]
Form the Hermitian block matrix with diagonal blocks $D_c-tI$ and
$-\Sigma^{-1}$ and off-diagonal block $\overline Q_c$.
Eliminating the lower block by a congruence gives the inertia of
$-\Sigma^{-1}$ plus that of $G_{\rm comp}-tI$. Eliminating the upper block
gives the inertia of $D_c-tI$ plus that of
$-\Sigma^{-1}-\overline Q_c^*(D_c-tI)^{-1}\overline Q_c$.
Equating negative counts proves the identity.
\end{proof}
The count measures compressed eigenvalues below $t$. Combined with
$\delta+\tau$, it brackets true thresholds using a matrix whose size is
the retained correction rank. Thresholds at poles require separate
treatment; near-zero numerical pivots or unresolved inertia do not give
a certificate. This representation does not imply that every plotted
residual forecast used a small secular solve: the full corrected forecast
may use the SVD of $\widehat J$, while the small inertia problem independently
checks selected spectral enclosures.

Gamma remains explicit in bulk attenuation and boundary factors. Increasing
$p$ improves the enclosure but increases correction rank. Report both:
small spectral error alone does not establish a compact explanatory theorem
if almost all dimensions have been retained.

\section{Target energy without isolating every eigenvector}
Eigenvalues alone do not determine which modes a target needs. Define
\[
S_K(s)=\frac{\|\mathbf1_{(-\infty,s]}(K)y\|^2}{\|y\|^2}.
\]
For a positive semidefinite kernel, this includes the nullspace for $s\ge0$.
Define $S_{\widehat K}$ similarly. A chosen buffer around the cutoff allows
a comparison without assuming a gap between every pair of eigenvalues.

\begin{theorem}[Buffered target-energy bounds]
For any $g>0$ and real $s$,
\[
\left(\sqrt{S_{\widehat K}(s-g)}-\frac\delta g\right)_+^2
\le S_K(s)
\le\min\!\left\{1,
\left(\sqrt{S_{\widehat K}(s+g)}+\frac\delta g\right)^2\right\}.
\]
\end{theorem}
\begin{proof}
Let $P=\mathbf1_{(-\infty,s]}(K)$ and
$Q_-=\mathbf1_{(-\infty,s-g]}(\widehat K)$.
On the corresponding restricted ranges, $X=(I-P)Q_-$ solves
\[
K_{>s}X-X\widehat K_{\le s-g}
=(I-P)(K-\widehat K)Q_-.
\]
The restricted spectra are separated by at least $g$. Shift both by $s-g$.
The integral solution uses a decaying exponential on the first range and
a nonexpanding exponential on the second, giving
\[
\|X\|\le\delta\int_0^\infty e^{-gt}\,dt=\delta/g.
\]
Therefore
\[
\|Q_-y\|\le\|Q_-Py\|+\|Q_-(I-P)y\|
\le\|Py\|+(\delta/g)\|y\|.
\]
This gives the lower bound. For the upper bound let
$Q_+=\mathbf1_{(-\infty,s+g]}(\widehat K)$.
The same separated-spectrum argument gives $\|P(I-Q_+)\|\le\delta/g$,
so
\[
\|Py\|\le\|Q_+y\|+(\delta/g)\|y\|.
\]
Normalize, square, and use $S_K(s)\le1$.
\end{proof}
If much target energy lies near the cutoff, a broad buffer can lose
information; if it is too narrow, $\delta/g$ can be large. This is an explicit
tradeoff, not an assumed eigengap. The theorem uses the corrected reference's
target projections. It does not claim a target-independent lower bound on
slow energy for all functions.

\section{From the spectral statement to acquisition time}
\begin{corollary}[Necessary times from slow target energy]
Suppose $\eta\|K\|\le1$, $0\le\eta s<1$, and $S_K(s)\ge a$. Then
\[
E(n)^2\ge a(1-\eta s)^{2n}.
\]
If $a>\varepsilon^2$ and $s>0$, reaching relative residual $\varepsilon$
requires
\[
n\ge\left\lceil
\frac{\log(\sqrt a/\varepsilon)}{-\log(1-\eta s)}
\right\rceil.
\]
Use Theorem~7's lower bound for $a$. For $s=0$, target energy in the
nullspace never decays; if $a>\varepsilon^2$, the threshold is unreachable.
For $a\le\varepsilon^2$, this cutoff gives no positive necessary time.
\end{corollary}
\begin{proof}
Expand the residual in the orthonormal eigenbasis. On eigenvalues at most
$s$, each squared contraction factor is at least $(1-\eta s)^{2n}$.
Sum the target energy in these modes and solve the inequality.
\end{proof}

\begin{corollary}[Full corrected forecasts and their error envelope]
Let
\[
\widehat E(n)=\frac{\|(I-\eta\widehat K)^ny\|}{\|y\|}.
\]
If $\eta(\|\widehat K\|+\delta)\le1$, both kernels have nonnegative
contraction factors and
\[
|E(n)-\widehat E(n)|\le n\eta\delta.
\]
\end{corollary}
\begin{proof}
The telescoping identity for powers writes their difference as $n$ products,
each containing one factor $\eta(K-\widehat K)$ and otherwise
norm-at-most-one factors. The operator difference is bounded by $n\eta\delta$.
Apply it to $y$ and use the reverse triangle inequality.
\end{proof}
Thus $\widehat E(n)+n\eta\delta\le\varepsilon$ certifies acquisition by $n$,
while $\widehat E(n)-n\eta\delta>\varepsilon$ certifies that it has not
occurred by $n$. The true residual is monotone under the stated step
condition. The bound need not remain useful at arbitrarily large $n$;
choose truncation for the time scale being tested.

With curvature normalization, $\eta=\kappa/\|K\|$ for $0<\kappa\le1$,
the step varies with gamma. Do not compare those runs as fixed-step runs.
If $\widehat L=\|\widehat K\|>\delta$, then
\[
\left\|\kappa\frac K{\|K\|}
-\kappa\frac{\widehat K}{\widehat L}\right\|
\le\frac{2\kappa\delta}{\widehat L-\delta}.
\]
Add and subtract $\kappa\widehat K/\|K\|$ and use
$|\|K\|-\widehat L|\le\delta$ to obtain this inequality.
The same telescoping proof uses this rate error instead of $\eta\delta$.
Alternatively, use the actual supplied step of an archived run directly.

These are exact-arithmetic statements with a perturbation allowance.
Numerical spectra or residuals close to machine precision need separate
numerical error bounds before their envelopes can be called certified.

\section{What scaling gamma with width actually preserves}
At a fixed grid-relative frequency $\theta=\omega h$,
\[
z=\frac{\pi|\theta|}{2\gamma h}
=\frac{\pi|\theta|}{2\lambda}.
\]
Keeping $\lambda=\gamma h$ fixed preserves smoothing attenuation at that
relative frequency. On a fixed interval, $h=L/N$, requiring gamma
proportional to $N$. If $\gamma h\to0$ and $\theta\ne0$ is fixed,
\[
M_\gamma(\theta/h)^2
=\frac{4z^2e^{-2z}}{(1-e^{-2z})^2}\longrightarrow0.
\]
This identifies an additional exponential smoothing penalty in the exact
periodic rates. It does not remove frequency dependence of $b_k$, raw width
factors, aliases, boundary corrections, or the step. It does not assert
width-independent learning times. For a fixed physical target frequency,
$\theta=\omega h$ shrinks as the grid is refined, so that is a different
scaling question.

\section{Numerical evaluation}
\label{sec:numerical-evaluation}
\subsection*{Protocol and what is computed}
The primary check uses the archived uniform geometry: $L=2$, $N=512$,
$q=16$, and $m=8193$. Its 512 core centers, terminal center at $L$, and
23 halo centers on each side give $W=559$
ordinary tanh features plus bias. Inputs include both endpoints. The four
common slopes are 8, 12, 16, and 64. Each saved step has
$\eta\|K\|\simeq0.5$. All raw readouts start at zero.

The target formulas below use the original centered coordinate
$x\in[-1,1]$; a translation gives the theorem's interval $[0,L]$.
The five targets are
$\sin(2\pi x)+\frac12\sin(6\pi x)+\frac14\sin(10\pi x)$,
$\exp(\sin(3\pi x))$, $(1+25x^2)^{-1}$, $\sqrt5x^2$, and
$\sqrt2\sin(2\pi x)$. The first is the primary target in the acquisition
figure. Two additional uniform geometries use $N=128,256$, the same
oversampling, and $\lceil\sqrt N\rceil$ halo centers on each side. This
gives 12 dictionaries and 60 target--dictionary cases. The added widths
have independent spectral references, not new executed GD trajectories.
No held-out generalization claim is made.

The bulk is synthesized from half-odd Fourier coefficients and their alias
blocks. Its Fourier tail is bounded below $10^{-20}$ per feature. For
each geometry and gamma, a finite scan over $B=1,\ldots,N/2-1$ and
$p=1,\ldots,256$ minimizes $2B+2p$ subject to an analytic kernel-tail
bound below $10^{-15}$. This selection uses the proved bounds, not original
eigenvalues or training observations. The primary cases use
$(B,p)=(23,29),(20,22),(15,22),(8,10)$.

After adding exact boundary columns, halo columns, bias, and the endpoint
row, we compute a rectangular SVD of $\widehat J$ for all target weights
and the residual curve. This is the SVD of the newly constructed model;
it is not the original tanh SVD. Separately, we construct the signed
coefficient-Gram correction, compress its eigenvalues below $10^{-12}$
in absolute value, and include discarded norm and reconstruction defect
in the spectral allowance. The compressed secular calculation independently
checks eigenvalue counts at normalized cutoffs $10^{-7},10^{-6},10^{-5}$,
matching all 36 reference counts across the 12 dictionaries. We do not
claim that every plotted forecast was calculated solely by the small
secular solve.

\begin{figure}[htbp]
\centering
\includegraphics[width=\linewidth]{uniform_grid_three_panel.png}
\caption{\textbf{Gamma's spectral mechanism carried through the actual finite
geometry.} A is the exact periodic multiplier-squared law, shown continuously
in frequency; it does not assert that ordinary-tanh modes are individual
waves. B compares the 26th normalized eigenvalue of the original tanh kernel
(circles) with the boundary-corrected prediction and its spectral enclosure.
C compares the primary target's executed GD acquisition times (circles)
with necessary--sufficient intervals from the constructed model. All three
panels use four slopes. B's enclosures are narrower than the plotted line;
their relative widths are quantified below. Numerical spectral bounds use
the disclosed arithmetic sensitivity allowance. The primary timing
intervals in C contain independently certified archived intervals.}
\end{figure}

\subsection*{What is now tight, and what remains a calculation}
Under the disclosed FP64 sensitivity allowance, the gamma-8 26th normalized eigenvalue is enclosed by
$[8.1417184,8.1418574]\times10^{-8}$, against the reference value
$8.1417879\times10^{-8}$. Its relative interval width is 0.00171\%.
At gamma 64, the interval is
$[4.978118981,4.978119066]\times10^{-4}$. Dividing these intervals bounds
the gamma-64/gamma-8 rate ratio by approximately $[6114.23,6114.34]$,
compared with the measured 6114.28. These compare the same ordered eigenvalue
at two slopes; they do not assert that its eigenvector remains unchanged.

The nominal kernel correction rank bounds are 306, 266, 246, and 170.
After the stated numerical compression, the retained signed ranks are
26, 26, 26, and 22 in a 560-dimensional coefficient space. Discarded
eigenvalues enter the error allowance. These reduced ranks are measured
properties, not a theorem that the correction always has rank 26. The
full analytic construction remains valid without compression.

For the primary target, the four necessary--sufficient intervals are
\[
\begin{gathered}
\left[15{,}767{,}081,15{,}829{,}763\right],\qquad
\left[186{,}056,186{,}059\right],\\
\left[61{,}792,61{,}793\right],\qquad
\left[16{,}013,16{,}013\right].
\end{gathered}
\]
They contain the executed hits $15{,}798{,}313$, $186{,}057$, $61{,}792$,
and $16{,}013$. The gamma-8 interval has width 0.397\% of the observed
time, and the predicted gamma-8/gamma-64 slowdown lies in $[984.64,988.56]$,
containing the observed 986.59.

All 60 cases have finite acquisition intervals enclosing their independent
spectral-reference hits. All 19 available archived GD hits are enclosed;
the quadratic/gamma-12 trajectory stopped at 200,000 updates, so its later
reference hit is not presented as an executed result. Across the sampled
residual curves of all 60 cases, the maximum discrepancy between constructed
and original models is $3.42\times10^{-14}$ in relative residual. This is
a numerical comparison, not a certificate for every update.

\subsection*{The target-energy bound is informative at the relevant rate}
At normalized cutoff $t=\eta s=10^{-6}$, the buffered theorem gives a
gamma-8 slow-energy fraction between 4.333177\% and 4.333388\%, containing
the reference 4.333330\%. Buffers are chosen by maximizing the lower and
minimizing the upper proved expressions over declared grids, without
consulting the true mass. The earlier simplified theorem gave a zero
lower bound at the same cutoff.

That lower energy bound alone gives a necessary time of 3,035,735 updates
to reach 1\% relative residual. Retaining every rate and target weight
gives the much tighter 15.8-million-update interval above. Thus the finite
theorem now quantifies a relevant slow component and its optimization
consequence. It does so through the corrected spectral model, not a
scalar gamma cap alone.

\begin{figure}[htbp]
\centering
\includegraphics[width=\linewidth]{uniform_grid_spectral_detail.png}
\caption{\textbf{Target alignment and limits of spectral resolution.}
A shows primary-target slow-energy fractions from the original tanh
spectrum, with buffered lower and upper envelopes from the constructed
model shaded. B shows spectral interval width divided by the original
eigenvalue, ordered from largest to smallest; only values exceeding the
allowance are plotted. Very small modes can have broad relative enclosures
even when relevant modes are sharply resolved. These are FP64 evaluations
of the proved formulas with an explicit arithmetic allowance, not independent
interval certificates for the spectra.}
\end{figure}

\subsection*{Arithmetic, independent checks, and reproducibility}
The exact-arithmetic kernel remainders in the primary construction are
all below $8\times10^{-16}$. Floating-point rounding dominates them.
We therefore add the feature-level sensitivity allowance
\[
e_{\rm fp}=64u\sqrt{W+1}(1+\gamma L+\log_2N)
+r_{\rm SVD}+\sigma_{\max}o_{\rm SVD},
\]
where $u$ is machine epsilon, $r_{\rm SVD}$ is the Frobenius reconstruction
residual, and $o_{\rm SVD}$ is the Frobenius orthogonality residual of left
singular vectors. This formula is a heuristic arithmetic allowance,
not a rigorous enclosure of transcendental evaluation and FFT rounding.
It is separate from the proved analytic remainder. Spectral intervals
also include the discarded signed-correction norm and measured Gram
reconstruction residual.

Multiplying $e_{\rm fp}$ by ten leaves all 60 cases with finite time
intervals. The primary gamma-8 interval widens to
$[15{,}495{,}367,16{,}123{,}116]$ and still resolves the large delay.
The four primary intervals above contain the existing independently
certified Arb endpoint intervals, so their validity for the nominal real
tanh problem also follows by monotonicity. This inherited endpoint
validation does not certify the new spectral numbers or heuristic allowance.

Independent checks covered 1,000 perturbed PSD examples for buffered
target mass, 500 signed secular-count examples, and 27 finite constructions
varying grid parity, oversampling, slope, and interval length, without
violations. Durable focused tests also compare iterative raw-readout GD
with the spectral solution and verify that neighboring differences
preserve the original parameter metric. The fixed-bandwidth check uses
$N=128,256,512$ and $\lambda=0.25,1,4$; the multiplier at $\theta=\pi/4$
is unchanged with width at fixed $\lambda$, as predicted.

The \href{\evidence/summary.json}{numerical record} stores every case,
source hashes, truncation choices, ranks, spectral intervals, target-energy
curves, and timing checks. Reproduce it from the repository root with
\begin{quote}\small\ttfamily
python -m experiments.expD36\_frozen\_gamma\_probe.uniform\_grid\_analysis
\end{quote}
Compile the note with
\begin{quote}\small\ttfamily
latexmk -pdf -outdir=/tmp/gamma-uniform-latex\\
\hspace*{1em}docs/gamma\_uniform\_grid\_theorem.tex
\end{quote}
Additional \href{\repo/docs/gamma_uniform_periodic_lemmas.md}{periodic-grid
proofs} accompany this note. No new GPU training was used.

\section{Scope of the result}
The periodic result is a sharp gamma-to-spectrum theorem. Its single-mode
timing consequences are exact, including inaccessible target components.
The finite ordinary-tanh result retains explicit attenuation but also
retains boundaries; its analytic remainder can be made arbitrarily small
at the cost of a larger correction. It supplies a structured spectral
problem and bounds transfer to the true kernel. It does not eliminate
that problem or guarantee a small correction rank in every regime.

All slopes are equal and frozen. Inputs and core centers are equally spaced,
with integer oversampling. Additional centers are allowed because their
columns are retained explicitly. Random sampling, nonuniform centers,
heterogeneous slopes, trained hidden parameters, and target-independent
monotonic acceleration are not established by these statements.
\end{document}
